%% Citas verificadas contra el PDF. Sin marcadores [VERIFY] pendientes.

\section{RESULTADOS Y DISCUSIÓN}

\subsection{Verificación preliminar del vector de características}

Para verificar que el vector de ocho características de la Tabla~\ref{tab:senales} contiene
información discriminante, se construyó un conjunto sintético de 10\,000 registros con tres
clases equilibradas. El conjunto se construyó ante la ausencia de un conjunto etiquetado
público de fichas de vendedores peruanos: sus variables y rangos se derivaron de los factores
de riesgo identificados en la Sección~II, y las etiquetas se asignaron mediante reglas
deterministas sobre esas mismas variables. Declarar esta procedencia es indispensable para
interpretar lo que sigue. Sus siete entradas corresponden a siete de las ocho de la
Tabla~\ref{tab:senales}. La veracidad de dominio, que este trabajo introduce, es la única que
no forma parte de ese vector. Se entrenaron cinco algoritmos de familias distintas sobre la
misma partición 80/20 con semilla fija, y se validaron además mediante validación cruzada
estratificada de cinco pliegues sobre el conjunto completo. Los resultados muestran que los
cinco algoritmos clasifican correctamente los 2\,000 registros de prueba, sin desviación
entre pliegues.

Los experimentos se ejecutaron sobre Ubuntu 24.04.3 LTS con Python 3.11 y scikit-learn, en un
procesador Intel Xeon E3-1240 v5 y sin aceleración gráfica. Las máquinas de soporte vectorial y
el clasificador de vecinos más cercanos se encadenaron a un normalizador de media cero y
varianza unitaria; los demás hiperparámetros se mantuvieron por defecto, con semilla fija,
porque el objetivo no era optimizar el desempeño sino comprobar si existía diferencia alguna
entre familias de algoritmos.

Los cinco algoritmos recuperan la regla generadora: que un clasificador de vecinos más
cercanos iguale a un ensamble de gradiente potenciado, y que ambos igualen a un árbol
individual, indica que la frontera de decisión es recuperable por cualquier método, de modo
que el problema es separable por construcción. Un árbol de decisión de profundidad cinco
reproduce la regla de etiquetado en su totalidad, y la importancia de las características se
concentra en dos: la puntuación del producto, con 0.51, y el volumen de ventas, con 0.46,
mientras que tres de las siete no contribuyen en absoluto.

El ejercicio vale como \textbf{verificación del canal de procesamiento y diagnóstico del
conjunto}, no como medida de eficacia. De ello se siguen dos consecuencias. La primera es que
estos datos \textbf{no permiten discriminar entre algoritmos}: la elección del bosque
aleatorio descrita en la Sección~III-B responde a requisitos de interpretabilidad y costo, no a
una comparación empírica que este conjunto no puede arbitrar. La segunda, más relevante, es
que las métricas obtenidas sobre él no estiman la calidad del nivel estimado de riesgo, sino
la fidelidad con que un modelo recupera la regla que generó sus propias etiquetas. La validez
del enfoque depende de que las características sean legibles en fichas reales y de que las
etiquetas provengan de una fuente independiente del vector de entrada, según establece el
protocolo de validación de la Sección~III-B.

\subsection{Amenazas a la validez}

Cuatro limitaciones condicionan lo que este diseño podrá concluir, y conviene explicitarlas
antes de su ejecución.

La primera afecta a la capa cliente. La extracción depende de selectores sobre el DOM, y se
ha advertido que las técnicas apoyadas en selectores predefinidos son sensibles a cambios de
maquetación \cite{ref24}. Una plataforma que rediseñe su ficha reducirá la cobertura de
características sin previo aviso, por lo que la Fase~1 debe repetirse periódicamente y no
tratarse como una medición única.

La segunda afecta al clasificador de riesgo. Sus etiquetas de entrenamiento se derivan
mediante reglas deterministas aplicadas a las mismas características que el modelo recibe como
entrada, de modo que el desempeño medido sobre ese conjunto refleja la recuperación de esa
regla y no el comportamiento frente a vendedores fraudulentos reales. La mitigación es previa
a cualquier medición de eficacia y condiciona el orden del trabajo futuro: construir un
conjunto cuyas etiquetas provengan de una fuente independiente del vector de características,
y hasta entonces no reportar métricas de desempeño del clasificador.

La tercera es de validez ecológica. Un experimento en sesión mide la capacidad de decisión en
condiciones de atención deliberada, mientras que la evidencia disponible sugiere que la
victimización ocurre justamente cuando esa atención falla \cite{ref27}. El diseño puede
establecer si la herramienta ayuda a quien la mira, no si protege a quien compra distraído.
Mitigarlo exige una segunda medición en uso real, con la extensión instalada durante semanas y
sobre las compras que el participante hubiera hecho de todos modos, que este protocolo no
cubre.

La cuarta es la más general: la arquitectura supone que el atacante no adapta su
comportamiento a las características observadas, suposición insostenible a mediano plazo
porque los defraudadores ajustan sus estrategias frente a los mecanismos de detección
\cite{ref28}. Como mitigación parcial, el registro de consultas descrito en la Sección~III-B
hace observable esa deriva y permite reentrenar cuando la distribución de las características
cambie; ninguna medición puntual, en cambio, conserva su validez indefinidamente.
